\subsection{A tree with one possibly large corner}

\begin{proposition}[Proved: a tree reduces to one bounded act cone]
\label{CP2:A-tree-one-large-corner-eleven}
Let $C$ be a finite connected category with at least two objects and acyclic support,
whose undirected comparability graph is a tree. Suppose
there is an object $v$ with $|\operatorname{End}(v)|\leq11$
and every other corner has order at most two. If $BC$ has
bounded rational homology, then all rational degree-two
cohomology products on $BC$ are zero.
\end{proposition}
\begin{proof}
A tree support has no chain
of three distinct objects, because transitivity would supply
the third edge of a triangle. All arrows incident with $v$
therefore point in the same direction.

Let $X$ be the full subcategory on $v$ and all its neighbors.
It is a nonempty source or sink star. Each component of
the graph outside this vertex set attaches to exactly one
vertex $w$ of $X$, and $w\ne v$. Let $Y$ be the full
subcategory on that component and its attachment vertex.
Every corner of $Y$ has order at most two. Its rational
corner algebras are semisimple, and its support has strict
height at most one. The relative bar complex therefore
gives $H_j(BY;\mathbb Q)=0$ for every $j\geq2$.
The same corners have rationally acyclic classifying spaces.

The nerves of $X$ and these $Y$ pieces cover $NC$ exactly:
each simplex uses at most two distinct objects and hence
lies over one vertex or one tree edge. Attach the pieces
one at a time. Each intersection is the corner category at
its attachment vertex $w$, which is rationally acyclic.
Mayer--Vietoris gives, in every degree $j\geq2$, the
restriction isomorphism
\[
 H^j(BC;\mathbb Q)\xrightarrow{\ \cong\ }H^j(BX;\mathbb Q).
\]
It is induced by the actual inclusion and respects products.
The corresponding homology sequences also show that $BX$
has bounded rational homology.

At every leaf of $X$, the trivial rational module is
projective on both sides, since its corner has order at
most two. Apply the actual orbit functors of
Lemma~\ref{CP2:A-high-orbits}, replacing these leaves by
trivial corners. They give rational cohomology-ring
isomorphisms and preserve bounded rational homology.
There is at least one leaf. Merge all the resulting trivial
leaves into one, retaining the disjoint union of their
stalks as the actual finite act of $\operatorname{End}(v)$.
There are no compositions between different leaves, so this
is a functor and its nerve realization identifies exactly
those finitely many leaf vertices. On a connected CW space,
identifying $r$ distinct points adds a wedge of $r-1$ circles
up to homotopy: this follows by attaching a contractible
tree joining them and then collapsing that tree. Thus the
quotient map is an isomorphism in cohomological degrees at
least two, preserving products, and the quotient still has
bounded rational homology.

The quotient is the ordinary act cone over the corner at
$v$, or its opposite. Theorem~\ref{CP2:A-bounded-eleven-acts}
kills every rational degree-two product there. Pullback
through the specified quotient and orbit functors kills
all such products on $BX$. The degree-four restriction
isomorphism then kills all products on $BC$.
\end{proof}
